% Chapter 18: Modern Combinatorics Breakthroughs
\let\LocalMacro\relax

\chapter{Modern Combinatorics: Cap Sets, Sensitivity, Sunflowers, and EFL}

\section{Prerequisites}
This chapter assumes familiarity with:
\begin{itemize}
\item Basic combinatorics (binomial coefficients, extremal combinatorics)
\item Polynomial method (Dvir's finite field Kakeya)
\item Probabilistic method (Lovász Local Lemma)
\item Basic theoretical computer science (boolean functions, decision trees)
\end{itemize}

\section{Introduction}
The years 2016--2022 saw an unprecedented convergence of techniques in combinatorics: the polynomial method, entropy arguments, and Fourier analysis on the hypercube resolved decades-old problems in rapid succession. This chapter surveys five landmark results.

\section{The Cap Set Problem}

\subsection{The Problem}
A \emph{cap set} in $\mathbb{F}_3^n$ is a subset with no three-term arithmetic progression (i.e., no distinct $x, y, z$ with $x+y+z=0$). Let $r_3(n)$ be the maximum size of a cap set in $\mathbb{F}_3^n$.

\begin{conjecture}[Frankl-Graham-Rödl, 1987]
$r_3(n) = o(3^n)$.
\end{conjecture}

The trivial bound is $3^n$. Meshulam (1995) proved $r_3(n) \le 2 \cdot 3^n / n$. The best pre-2016 bound was $r_3(n) \le 3^n / n^{1+c}$.

\subsection{Ellenberg-Gijswijt (2016)}

\begin{theorem}[Ellenberg-Gijswijt \cite{ellenberggijswijt2017}]
$r_3(n) \le 2.756^n$ for sufficiently large $n$.
\end{theorem}

Their proof used the \emph{slice rank} method (a variant of the polynomial method):
\begin{enumerate}
\item Represent the cap set condition by a tensor $T: \mathbb{F}_3^n \times \mathbb{F}_3^n \times \mathbb{F}_3^n \to \{0,1\}$.
\item Show that the slice rank of $T$ is at most $3 \cdot M$, where $M$ is the number of monomials $x_1^{e_1} \cdots x_n^{e_n}$ with $e_i \in \{0,1,2\}$ and $\sum e_i < 2n/3$.
\item Use Stirling's formula to bound $M \le 2.756^n$.
\end{enumerate}

\begin{highlightbox}
The slice rank method is a powerful generalization of the polynomial method: instead of vanishing of a polynomial, one uses the rank of a tensor. It has since been applied to many problems (sunflower, Erdős-Ginzburg-Ziv, etc.).
\end{highlightbox}

\section{The Sunflower Conjecture}

\subsection{The Problem}
A \emph{sunflower} with $k$ petals is a family of $k$ sets whose pairwise intersections are all equal to a common core $C$. The Erdős-Rado sunflower conjecture (1960):

\begin{conjecture}
Any family of sets of size $w$ with more than $c(w)^k$ sets contains a sunflower with $k$ petals.
\end{conjecture}

Erdős offered \$1000 for a proof. The best bound was $(w!)^{k-1}$.

\subsection{Alweiss, Lovett, Wu, Zhang (2022)}

\begin{theorem}[ALWZ \cite{alweiss2022}]
Any family of sets of size $w$ with more than $(C \log w)^{k-1} w!$ sets contains a sunflower with $k$ petals.
\end{theorem}

Their proof uses:
\begin{enumerate}
\item \emph{Spread approximations}: A set family is ``spread'' if no element appears too often. Any family can be approximated by a spread family.
\item \emph{Random sampling}: A random subset of a spread family contains a sunflower with high probability.
\item \emph{Entropy arguments}: Use the Lopsided Lovász Local Lemma to find the sunflower.
\end{enumerate}

\section{The Sensitivity Conjecture}

\subsection{The Problem}
For a boolean function $f: \{0,1\}^n \to \{0,1\}$, the \emph{sensitivity} $s(f)$ is the maximum number of coordinates whose flip changes $f$. The \emph{block sensitivity} $bs(f)$ allows flipping blocks. The \emph{degree} $\deg(f)$ is the degree of the polynomial representing $f$ over $\mathbb{R}$.

\begin{conjecture}[Sensitivity Conjecture, Nisan-Szegedy 1992]
$s(f)$ and $bs(f)$ are polynomially related: $bs(f) \le s(f)^c$ for some constant $c$.
\end{conjecture}

It was known that $\deg(f) \le s(f)^2$ and $bs(f) \le \deg(f)^2$, so the conjecture would follow from $s(f) \le \deg(f)^c$.

\subsection{Huang (2019)}

\begin{theorem}[Huang \cite{huang2019}]
$s(f) \ge \sqrt{\deg(f)}$, hence $bs(f) \le s(f)^4$.
\end{theorem}

Huang's proof:
\begin{enumerate}
\item Consider the $n$-dimensional hypercube graph $Q_n$.
\item Construct a signing of the adjacency matrix $A$ such that $A^2 = n I$.
\item By Cauchy interlacing, the subgraph induced by $f^{-1}(1)$ has an eigenvalue $\ge \sqrt{n}$.
\item This implies sensitivity $\ge \sqrt{n}$.
\end{enumerate}

\begin{highlightbox}
Huang's proof is famously short (2 pages) and uses only spectral graph theory. It resolved a 27-year-old problem that had resisted Fourier analysis, probability, and complexity theory.
\end{highlightbox}

\section{The Kahn-Kalai Conjecture (Expectation-Threshold)}

\subsection{The Problem}
For an increasing property $\mathcal{P}$ on subsets of $[n]$, the \emph{threshold} $p_c(\mathcal{P})$ is the $p$ where $\mathbb{P}(G(n,p) \in \mathcal{P}) = 1/2$. The \emph{expectation threshold} $q(\mathcal{P})$ is the smallest $p$ where the expected number of witnesses is 1.

\begin{conjecture}[Kahn-Kalai, 2006]
$p_c(\mathcal{P}) \le C \log \ell \cdot q(\mathcal{P})$, where $\ell$ is the size of the smallest witness.
\end{conjecture}

\subsection{Park-Pham (2022)}

\begin{theorem}[Park-Pham \cite{parkpham2022}]
The Kahn-Kalai conjecture is true: $p_c(\mathcal{P}) \le C \log \ell \cdot q(\mathcal{P})$.
\end{theorem}

Their proof uses:
\begin{enumerate}
\item \emph{Spread approximations}: Any family can be approximated by a spread family.
\item \emph{Weighted random sampling}: A clever two-stage random process to find the threshold.
\item \emph{Entropy compression}: To bound the failure probability.
\end{enumerate}

\section{The Erdős-Faber-Lovász Conjecture}

\subsection{The Problem}
\begin{conjecture}[Erdős-Faber-Lovász, 1972]
If $G$ is the union of $n$ cliques of size $n$, each pair intersecting in at most one vertex, then $\chi(G) = n$.
\end{conjecture}

Erdős offered \$500. The conjecture was open for 49 years.

\subsection{Kang, Kelly, Kühn, Methuku, Osthus (2022)}

\begin{theorem}[KKKMO \cite{kang2022}]
EFL conjecture is true for sufficiently large $n$.
\end{theorem}

Their proof uses:
\begin{enumerate}
\item \emph{Absorption method}: Build the coloring iteratively, absorbing leftover vertices at the end.
\item \emph{Iterative absorption}: Absorb small configurations step by step.
\item \emph{Robust expanders}: Use expansion properties of the hypergraph to guarantee absorption works.
\end{enumerate}

\section{Impact}
\begin{itemize}
\item \textbf{Polynomial method 2.0}: Slice rank, spread approximations, and entropy arguments form a new toolkit.
\item \textbf{Connections to TCS}: Sensitivity, threshold phenomena, and boolean function analysis are now deeply connected to combinatorics.
\item \textbf{New conjectures}: These methods suggest new bounds for Ramsey theory, graph coloring, and extremal set theory.
\end{itemize}

\section{Exercises}

\begin{exercise}[Basic]
Show that a cap set in $\mathbb{F}_3^n$ has no three points on a line. Compute the maximum cap set size for $n=2$.
\end{exercise}

\begin{exercise}[Intermediate]
Explain how the slice rank of a tensor generalizes the rank of a matrix. Why does slice rank bound cap sets?
\end{exercise}

\begin{exercise}[Challenge]
Sketch Huang's proof of the sensitivity conjecture using the signed adjacency matrix of the hypercube.
\end{exercise}

\section{Timeline}

\begin{tabularx}{\linewidth}{lX}
\toprule
\textbf{Year} & \textbf{Event} \\ \midrule
2016 & Ellenberg-Gijswijt solve cap set with slice rank \\ 
2019 & Huang proves sensitivity conjecture \\ 
2022 & ALWZ prove sunflower conjecture \\ 
2022 & Park-Pham prove Kahn-Kalai conjecture \\ 
2022 & KKKMO prove Erdős-Faber-Lovász conjecture \\ 
2022 & Huh wins Fields Medal (log-concavity in combinatorics) \\
\bottomrule
\end{tabularx}

\section{Citations}

\begin{enumerate}
\item Ellenberg, J., \& Gijswijt, D. (2017). ``On large subsets of $\mathbb{F}_3^n$ with no three-term arithmetic progression.'' Annals of Mathematics, 185(1), 339--343.
\item Tao, T. (2016). ``A symmetric formulation of the Croot-Lev-Pach-Ellenberg-Gijswijt capset bound.'' Blog post.
\item Huang, H. (2019). ``Induced subgraphs of hypercubes and a proof of the sensitivity conjecture.'' Annals of Mathematics, 190(3), 949--955.
\item Alweiss, R., Lovett, S., Wu, K., \& Zhang, J. (2022). ``Improved bounds for the sunflower lemma.'' Annals of Mathematics, 196(3), 977--1021.
\item Park, J., \& Pham, H. T. (2022). ``A proof of the Kahn-Kalai conjecture.'' arXiv:2203.17207.
\item Kang, D., Kelly, T., Kühn, D., Methuku, A., \& Osthus, D. (2022). ``A proof of the Erdős-Faber-Lovász conjecture.'' arXiv:2101.04686.
\end{enumerate}
